% =============================================================================
% Section 1: Introduction
% =============================================================================

\section{Introduction}
\label{sec:introduction}

Computational fluid dynamics (CFD) has matured into an indispensable tool 
across engineering, atmospheric science, and fundamental physics research. 
However, a persistent gap remains between the mathematical formulations 
governing fluid motion and their numerical implementations. Discretization 
errors, boundary condition artifacts, and algorithmic approximations accumulate 
silently, undermining confidence in simulation results—particularly for 
safety-critical applications in aerospace (DO-178C Level~A) and biomedical 
engineering (FDA Class~III).

\subsection{Motivation}

The LeanFlow project addresses this gap through a \emph{dual-scale} 
approach: coupling a pseudo-spectral Navier-Stokes solver with the Lean~4 
interactive theorem prover to create a pipeline where:
\begin{enumerate}
  \item The continuum PDE semantics are axiomatized in Lean~4 (Section~\ref{sec:formalization}).
  \item The numerical solver produces empirical results validated against 
        published reference solutions (Section~\ref{sec:results}).
  \item Formal specification roadmaps bridge the mathematical theory to the 
        implementation, enabling incremental proof construction (Section~\ref{sec:formalization}).
\end{enumerate}

\subsection{Contributions}

This paper makes the following contributions:
\begin{itemize}
  \item \textbf{Reproducible benchmark execution} of 10 canonical PDE use cases 
        (UC7--UC16) with wall-clock timing, L$_2$ error metrics, and SHA-256 
        certification seals (Section~\ref{sec:results}).
  \item \textbf{Lean~4 axiomatic formalization} of the incompressible 
        Navier-Stokes equations, Boussinesq approximation, inviscid Euler 
        limit, and spectral enstrophy bounds (Section~\ref{sec:formalization}).
  \item \textbf{Quantitative comparison} against established reference solutions including the PDEBench dataset~\citep{takamoto2022pdebench}, the lid-driven cavity profiles of \citet{ghia1982high}, Dedalus~\citep{burns2020dedalus}, Athena++~\citep{stone2020athena}, and the JHTDB turbulence database~\citep{li2008public} (Section~\ref{sec:comparison}).
  \item \textbf{Reproducible protocol} with deterministic random seeds, 
        certified JSON output, and SHA-256 result sealing.
\end{itemize}

\subsection{Organization}

Section~\ref{sec:methodology} describes the solver architecture and benchmark 
protocol. Section~\ref{sec:formalization} presents the Lean~4 formal 
specification roadmaps. Section~\ref{sec:results} reports the experimental 
results. Section~\ref{sec:comparison} compares against traditional solvers. 
Section~\ref{sec:conclusion} discusses implications and future directions.
